\documentclass[]{edoxws}

\renewcommand\examname{2. Arbeit }
\renewcommand\sign{B. Falkner}

\begin{document}

\pdate{März 2026}
\begin{exam}{Informatik, WP 09}

% ==================================================
\exercise[Grundlagen – Begriffe der Verschlüsselung €]{

  \subexercise{
    Kreuze alle richtigen Aussagen an:
    \begin{itemize}
      \item[$\Box$] Klartext ist der ursprüngliche, lesbare Text.
      \item[$\Box$] Chiffretext ist eine unverschlüsselte Nachricht.
      \item[$\Box$] Ein Schlüssel bestimmt, wie verschlüsselt wird.
      \item[$\Box$] Verschlüsselung dient dem Schutz von Daten.
    \end{itemize}

    \begin{enum}
      \item Klartext ist der ursprüngliche, lesbare Text.
      \item Chiffretext ist eine unverschlüsselte Nachricht.
      \item Ein Schlüssel bestimmt, wie verschlüsselt wird.
      \item Verschlüsselung dient dem Schutz von Daten.
    \end{enum}

    \solve{
      \begin{itemize}
        \item[$\boxtimes$] Klartext ist der ursprüngliche, lesbare Text.\pp
        \item[$\Box$] Chiffretext ist eine unverschlüsselte Nachricht.
        \item[$\boxtimes$] Ein Schlüssel bestimmt, wie verschlüsselt wird.\pp
        \item[$\boxtimes$] Verschlüsselung dient dem Schutz von Daten.\pp
      \end{itemize}
      \pp[1]
    }
  }

  \subexercise{
    Welche Begriffe gehören zur Verschlüsselung?
    \begin{itemize}
      \item[$\Box$] Klartext
      \item[$\Box$] Schlüssel
      \item[$\Box$] Schleife
      \item[$\Box$] Chiffretext
      \item[$\Box$] Bedingung
    \end{itemize}

    \solve{
      \begin{itemize}
        \item[$\boxtimes$] Klartext\pp
        \item[$\boxtimes$] Schlüssel\pp
        \item[$\Box$] Schleife
        \item[$\boxtimes$] Chiffretext\pp
        \item[$\Box$] Bedingung
      \end{itemize}
      \pp[2]
    }
  }
}

% ==================================================
\exercise[Caesar-Verschlüsselung]{

  \subexercise{
    Kreuze an:
    \begin{itemize}
      \item[$\Box$] Bei der Caesar-Verschlüsselung werden Buchstaben verschoben.
      \item[$\Box$] Es gibt unendlich viele mögliche Schlüssel.
      \item[$\Box$] Häufigkeitsanalyse kann das Verfahren angreifen.
      \item[$\Box$] Die Caesar-Verschlüsselung gilt heute als sicher.
    \end{itemize}

    \solve{
      \begin{itemize}
        \item[$\boxtimes$] Bei der Caesar-Verschlüsselung werden Buchstaben verschoben.\pp
        \item[$\Box$] Es gibt unendlich viele mögliche Schlüssel.
        \item[$\boxtimes$] Häufigkeitsanalyse kann das Verfahren angreifen.\pp
        \item[$\Box$] Die Caesar-Verschlüsselung gilt heute als sicher.
      \end{itemize}
      \pp[2]
    }
  }

  \subexercise{
    Warum ist die Caesar-Verschlüsselung unsicher?
    \begin{itemize}
      \item[$\Box$] Es gibt nur wenige mögliche Schlüssel.
      \item[$\Box$] Jeder kennt das Alphabet.
      \item[$\Box$] Häufige Buchstaben können erkannt werden.
      \item[$\Box$] Der Schlüssel ist immer öffentlich.
    \end{itemize}

    \solve{
      \begin{itemize}
        \item[$\boxtimes$] Es gibt nur wenige mögliche Schlüssel.\pp
        \item[$\Box$] Jeder kennt das Alphabet.
        \item[$\boxtimes$] Häufige Buchstaben können erkannt werden.\pp
        \item[$\Box$] Der Schlüssel ist immer öffentlich.
      \end{itemize}
      \pp[2]
    }
  }

  \subexercise{
    Häufigkeitsanalyse\\
    Ein Angreifer untersucht, wie oft einzelne Buchstaben im Chiffretext vorkommen.

    Kreuze an, welche Aussagen zur Häufigkeitsanalyse passen:
    \begin{itemize}
      \item[$\Box$] Im Deutschen kommt der Buchstabe E besonders häufig vor.
      \item[$\Box$] Häufigkeitsanalyse funktioniert nur bei sehr kurzen Texten.
      \item[$\Box$] Durch häufige Buchstaben kann der Schlüssel vermutet werden.
      \item[$\Box$] Häufigkeitsanalyse ist nur bei modernen Verfahren möglich.
    \end{itemize}

    \solve{
      \begin{itemize}
        \item[$\boxtimes$] Im Deutschen kommt der Buchstabe E besonders häufig vor.\pp
        \item[$\Box$] Häufigkeitsanalyse funktioniert nur bei sehr kurzen Texten.
        \item[$\boxtimes$] Durch häufige Buchstaben kann der Schlüssel vermutet werden.\pp
        \item[$\Box$] Häufigkeitsanalyse ist nur bei modernen Verfahren möglich.
      \end{itemize}
      \pp[2]
    }
  }
}

% ==================================================
\exercise[Symmetrische Verschlüsselung]{

  \subexercise{
    Kreuze alle richtigen Aussagen an:
    \begin{itemize}
      \item[$\Box$] Sender und Empfänger benutzen denselben Schlüssel.
      \item[$\Box$] Der Schlüssel muss geheim bleiben.
      \item[$\Box$] Es gibt einen öffentlichen und einen privaten Schlüssel.
      \item[$\Box$] Der Schlüsselaustausch ist ein Sicherheitsproblem.
    \end{itemize}

    \solve{
      \begin{itemize}
        \item[$\boxtimes$] Sender und Empfänger benutzen denselben Schlüssel.\pp
        \item[$\boxtimes$] Der Schlüssel muss geheim bleiben.\pp
        \item[$\Box$] Es gibt einen öffentlichen und einen privaten Schlüssel.
        \item[$\boxtimes$] Der Schlüsselaustausch ist ein Sicherheitsproblem.\pp
      \end{itemize}
      \pp[1]
    }
  }

  \subexercise{
    Welche Beispiele nutzen symmetrische Verschlüsselung?
    \begin{itemize}
      \item[$\Box$] Passwörter
      \item[$\Box$] Verschlüsselte Datei auf USB-Stick
      \item[$\Box$] HTTPS beim ersten Seitenaufruf
      \item[$\Box$] Öffentlicher Schlüssel
    \end{itemize}

    \solve{
      \begin{itemize}
        \item[$\boxtimes$] Passwörter\pp
        \item[$\boxtimes$] Verschlüsselte Datei auf USB-Stick\pp
        \item[$\Box$] HTTPS beim ersten Seitenaufruf
        \item[$\Box$] Öffentlicher Schlüssel
      \end{itemize}
      \pp[2]
    }
  }

  \subexercise{
    Vorteile und Nachteile symmetrischer Verschlüsselung\\
    Kreuze an, welche Aussagen richtig sind:

    \begin{itemize}
      \item[$\Box$] Symmetrische Verschlüsselung ist sehr schnell.
      \item[$\Box$] Symmetrische Verschlüsselung eignet sich gut für große Datenmengen.
      \item[$\Box$] Der geheime Schlüssel muss sicher übertragen werden.
      \item[$\Box$] Die Schlüsselübertragung ist unproblematisch.
      \item[$\Box$] Symmetrische Verfahren sind langsamer als asymmetrische.
    \end{itemize}

    \solve{
      \begin{itemize}
        \item[$\boxtimes$] Symmetrische Verschlüsselung ist sehr schnell.\pp
        \item[$\boxtimes$] Symmetrische Verschlüsselung eignet sich gut für große Datenmengen.\pp
        \item[$\boxtimes$] Der geheime Schlüssel muss sicher übertragen werden.\pp
        \item[$\Box$] Die Schlüsselübertragung ist unproblematisch.
        \item[$\Box$] Symmetrische Verfahren sind langsamer als asymmetrische.
      \end{itemize}
      \pp[2]
    }
  }
}

% ==================================================
\exercise[Asymmetrische Verschlüsselung / Public-Key]{

  \subexercise{
    Kreuze an:
    \begin{itemize}
      \item[$\Box$] Jeder besitzt einen öffentlichen und einen privaten Schlüssel.
      \item[$\Box$] Der private Schlüssel darf weitergegeben werden.
      \item[$\Box$] Mit dem öffentlichen Schlüssel wird verschlüsselt.
      \item[$\Box$] Mit dem privaten Schlüssel wird entschlüsselt.
      \item[$\Box$] Es ist kein geheimer Schlüsselaustausch nötig.
    \end{itemize}

    \solve{
      \begin{itemize}
        \item[$\boxtimes$] Jeder besitzt einen öffentlichen und einen privaten Schlüssel.\pp
        \item[$\Box$] Der private Schlüssel darf weitergegeben werden.
        \item[$\boxtimes$] Mit dem öffentlichen Schlüssel wird verschlüsselt.\pp
        \item[$\boxtimes$] Mit dem privaten Schlüssel wird entschlüsselt.\pp
        \item[$\boxtimes$] Es ist kein geheimer Schlüsselaustausch nötig.\pp
      \end{itemize}
      \pp[1]
    }
  }

  \subexercise{
    Ordne zu: Welche Aussagen passen zur asymmetrischen Verschlüsselung?
    \begin{itemize}
      \item[$\Box$] sehr schnell
      \item[$\Box$] zwei verschiedene Schlüssel
      \item[$\Box$] sicherer Schlüsselaustausch
      \item[$\Box$] gleicher Schlüssel für alle
    \end{itemize}

    \solve{
      \begin{itemize}
        \item[$\Box$] sehr schnell
        \item[$\boxtimes$] zwei verschiedene Schlüssel\pp
        \item[$\boxtimes$] sicherer Schlüsselaustausch\pp
        \item[$\Box$] gleicher Schlüssel für alle
      \end{itemize}
      \pp[2]
    }
  }

  \subexercise{
  Erkläre in eigenen Worten Schritt für Schritt, wie eine asymmetrische
  Verschlüsselung funktioniert.  
  Gehe dabei auf den öffentlichen und den privaten Schlüssel ein.

  \medskip
  \vspace{0.6cm}
  \hrule
  \vspace{0.6cm}
  \hrule
  \vspace{0.6cm}
  \hrule
  \vspace{0.6cm}
  \hrule
  \vspace{0.6cm}
  \hrule
  \vspace{0.6cm}
  \hrule
  \vspace{0.6cm}
  \hrule
  \vspace{0.6cm}
  \hrule

  \solve{
    Zuerst besitzt der Empfänger einen öffentlichen und einen privaten Schlüssel\pp.
    Der öffentliche Schlüssel wird weitergegeben\pp.
    Der Sender verschlüsselt die Nachricht mit dem öffentlichen Schlüssel des Empfängers\pp.
    Die verschlüsselte Nachricht wird übertragen\pp.
    Nur der private Schlüssel des Empfängers kann die Nachricht entschlüsseln\pp.
    Dadurch kann kein Angreifer den Inhalt lesen\pp.
  }
}


}




\exercise[Verschlüsselung beim Online-Banking]{

\vspace{0.5cm}
\textbf{
  Wie funktioniert Verschlüsselung beim Online-Banking?}

Wenn du dich beim Online-Banking anmeldest, müssen sehr sensible Daten geschützt werden,
zum Beispiel deine Zugangsdaten oder Überweisungen..

Zu Beginn der Verbindung baut dein Browser eine \textbf{HTTPS-Verbindung} zur Bank auf.
Dabei nutzt die Bank ein \textbf{Zertifikat}, das ihren öffentlichen Schlüssel enthält.
So kann dein Browser sicher überprüfen, ob er wirklich mit der echten Bank spricht
und nicht mit einem Betrüger.

Anschließend wird über diese sichere Verbindung ein \textbf{geheimer Sitzungsschlüssel}
ausgetauscht. Ab diesem Moment erfolgt die gesamte weitere Kommunikation
mit diesem Sitzungsschlüssel.
Das hat einen wichtigen Grund: so ist es deutlich
\textbf{schneller} und eignet sich besser für größere Datenmengen.

Während der Anmeldung werden Benutzername, Passwort und weitere Sicherheitsmerkmale
(z.\,B. TAN) über die verschlüsselte Verbindung übertragen.
So werden die Schutzziele der Datensicherheit erfüllt:
\textbf{Vertraulichkeit} (niemand kann mitlesen),
\textbf{Integrität} (Daten werden nicht verändert)
und \textbf{Authentizität} (die Bank ist echt).

\vspace{0.5cm}

% --------------------------------------------------
\subexercise{
\textbf{Ablauf der Verschlüsselung:} Kreuze alle richtigen Aussagen an:

\begin{itemize}
  \item[$\Box$] Beim Online-Banking wird zuerst asymmetrisch verschlüsselt.
  \item[$\Box$] Der öffentliche Schlüssel der Bank ist geheim.
  \item[$\Box$] Asymmetrische Verschlüsselung dient dem sicheren Schlüsselaustausch.
  \item[$\Box$] Symmetrische Verschlüsselung ist schneller als asymmetrische.
  \item[$\Box$] Nach dem Verbindungsaufbau wird die gesamte Kommunikation symmetrisch verschlüsselt.
\end{itemize}

\solve{
\begin{itemize}
  \item[$\boxtimes$] Beim Online-Banking wird zuerst asymmetrisch verschlüsselt.\pp
  \item[$\Box$] Der öffentliche Schlüssel der Bank ist geheim.
  \item[$\boxtimes$] Asymmetrische Verschlüsselung dient dem sicheren Schlüsselaustausch.\pp
  \item[$\boxtimes$] Symmetrische Verschlüsselung ist schneller als asymmetrische.\pp
  \item[$\boxtimes$] Nach dem Verbindungsaufbau wird die gesamte Kommunikation symmetrisch verschlüsselt.\pp
\end{itemize}
\pp[1]
}
}

% --------------------------------------------------
\subexercise{
\textbf{Vorteile der Kombination:} Warum werden beim Online-Banking beide Verfahren kombiniert?

\begin{itemize}
  \item[$\Box$] Asymmetrische Verschlüsselung ist besonders schnell.
  \item[$\Box$] Symmetrische Verschlüsselung eignet sich für große Datenmengen.
  \item[$\Box$] Der sichere Austausch eines geheimen Schlüssels wird ermöglicht.
  \item[$\Box$] Dadurch kann vollständig auf Verschlüsselung verzichtet werden.
\end{itemize}

\solve{
\begin{itemize}
  \item[$\Box$] Asymmetrische Verschlüsselung ist besonders schnell.
  \item[$\boxtimes$] Symmetrische Verschlüsselung eignet sich für große Datenmengen.\pp
  \item[$\boxtimes$] Der sichere Austausch eines geheimen Schlüssels wird ermöglicht.\pp
  \item[$\Box$] Dadurch kann vollständig auf Verschlüsselung verzichtet werden.
\end{itemize}
\pp[2]
}
}

% --------------------------------------------------
\subexercise{
Beschreibe in eigenen Worten den \textbf{kompletten Ablauf der Anmeldung}
beim Online-Banking.
Gehe dabei auf folgende Punkte ein:
\begin{itemize}
  \item Aufbau der sicheren Verbindung
  \item Rolle der Verschlüsselung zu Beginn
  \item Übergang zur dauerhaften Verschlüsselung der Verbindung 
  \item Übertragung der Zugangsdaten
\end{itemize}

\medskip
\hrule
\vspace{0.6cm}
\hrule
\vspace{0.6cm}
\hrule
\vspace{0.6cm}
\hrule
\vspace{0.6cm}
\hrule
\vspace{0.6cm}
\hrule
\vspace{0.6cm}
\hrule
\vspace{0.6cm}
\hrule

\solve{
Beim Aufruf der Bankseite wird zuerst eine HTTPS-Verbindung aufgebaut\pp.
Dabei überprüft der Browser mithilfe des Zertifikats den öffentlichen Schlüssel der Bank\pp.
Anschließend wird asymmetrische Verschlüsselung genutzt, um einen geheimen Sitzungsschlüssel auszutauschen\pp.
Danach erfolgt die weitere Kommunikation symmetrisch, da dies schneller ist\pp.
So sind die Daten vertraulich, unverändert und der Absender ist eindeutig\pp.
}
}

}

% ==================================================
\exercise[Schutzziele der Datensicherheit]{

  \subexercise{
    Kreuze die Schutzziele an:
    \begin{itemize}
      \item[$\Box$] Vertraulichkeit
      \item[$\Box$] Integrität
      \item[$\Box$] Geschwindigkeit
      \item[$\Box$] Authentizität
      \item[$\Box$] Verbindlichkeit
    \end{itemize}

    \solve{
      \begin{itemize}
        \item[$\boxtimes$] Vertraulichkeit\pp
        \item[$\boxtimes$] Integrität\pp
        \item[$\Box$] Geschwindigkeit
        \item[$\boxtimes$] Authentizität\pp
        \item[$\boxtimes$] Verbindlichkeit\pp
      \end{itemize}
      \pp[1]
    }
  }

  \subexercise{
    Welche Aussage beschreibt Integrität?
    \begin{itemize}
      \item[$\Box$] Nur Berechtigte dürfen lesen.
      \item[$\Box$] Daten wurden nicht verändert.
      \item[$\Box$] Der Absender ist bekannt.
      \item[$\Box$] Niemand kann etwas abstreiten.
    \end{itemize}

    \solve{
      Daten wurden nicht verändert\pp[4]
    }
  }
}

% ==================================================
\exercise[Ordnungspunkte]{
  \subexercise*{
    SEB, sprachlich ordentlich, Form und Inhalt sind gut strukturiert und verständlich.
    \solve{
      \begin{itemize}
        \item SEB \pp
        \item Sprachlich ordentlich \pp
        \item Form und Inhalt gut strukturiert \pp
      \end{itemize}
    }
  }
}

\end{exam}
\end{document}